%!TEX root = ./main.tex

\section[Angle]{Ruler 1: Angle, Read Across the Antennas}

% SECTION TIMING: 22:00--40:00 (6 frames). This is ArrayTrack's story, told on the
% running example. The section's one sentence: the SAME wavefront hits each antenna
% at a slightly different time, so the phase STEP between neighbors encodes the
% direction -- and a step is a difference, so wrapping is beaten.

% Teaching note (150 s): pure intuition, zero math. Snap your fingers on the left of
% the room: which ear heard it first? The brain does AoA from micro-delays; the AP
% does it from phase. Only then advance to the geometry.
\begin{frame}{Two ears, one voice}
  \vspace{0.3em}
  \centering
  \begin{tikzpicture}[scale=0.9,transform shape]
    % incoming plane wavefronts at 30 degrees from boresight (i.e. tilted)
    \foreach \s in {0,0.55,1.1}{
      \draw[softgreen!70,thick] ($(1.1,3.3)+(-60:\s)+(150:2.2)$)
        -- ($(1.1,3.3)+(-60:\s)-(150:2.2)$);}
    \draw[-{Latex[length=2.6mm]},very thick,softgreen]
      (1.1,3.3) ++(150:1.6) ++(-60:-0.7) -- ++(-60:1.6);
    \node[font=\scriptsize\bfseries,text=softgreen] at (3.4,3.6)
      {wavefronts from the phone, $\theta = 30^\circ$};
    % two antennas
    \anttri{(1.5,0.6)}{ranblue}
    \anttri{(3.5,0.6)}{ranblue}
    \node[font=\scriptsize,text=ranblue] at (1.5,0.1) {antenna 1};
    \node[font=\scriptsize,text=ranblue] at (3.5,0.1) {antenna 2};
    \draw[softgray,{Latex[length=1.6mm]}-{Latex[length=1.6mm]}]
      (1.5,-0.25) -- node[below,font=\scriptsize]{$d$} (3.5,-0.25);
    % extra path to antenna 2
    \draw[very thick,coreorange] (3.5,0.75) -- ++(120:1.0);
    \node[font=\scriptsize\bfseries,text=coreorange,align=left] at (5.6,1.5)
      {extra travel to antenna 2:\\$\Delta d = d\,\sin\theta$};
  \end{tikzpicture}

  \vspace{0.5em}
  \qa{The wavefront is a straight line. Why does antenna 2 hear it late?}
     {Because the line is \emph{tilted} by the arrival angle $\theta$. The lag
      \textbf{is} the direction --- your ears run this algorithm all day.}
\end{frame}

% Teaching note (210 s): the worked number is the point of the frame. d = 3 cm,
% lambda = 6 cm, theta = 30 deg: a quarter turn per antenna. Have them compute the
% step for theta = 0 (zero) and theta = 90 (half turn) -- the full dial.
\begin{frame}{The phase difference is the angle}
  \vspace{0.4em}
  \centering
  {\large $\Delta\varphi \;=\; \dfrac{2\pi}{\lambda}\, d \sin\theta$\par}
  \vspace{0.35em}
  {\small\textcolor{softgray}{extra path $d\sin\theta$, charged at one turn per
    wavelength}\par}

  \vspace{0.6em}
  \begin{tcolorbox}[colback=green!5,colframe=softgreen,boxrule=0.7pt,arc=2pt,
    left=5pt,right=5pt,top=4pt,bottom=4pt,width=0.9\textwidth]
    \small\centering
    Running example: $\lambda = 6$ cm, $d = 3$ cm, phone at $\theta = 30^\circ$
    \\[3pt]
    $\Delta\varphi = \frac{2\pi}{6\,\text{cm}} \cdot 3\,\text{cm} \cdot
     \sin 30^\circ = \pi/2$ \;---\; \textbf{a quarter turn per antenna}
  \end{tcolorbox}

  \vspace{0.5em}
  \qa{Why do arrays space antennas exactly $d = \lambda/2$?}
     {Wrapping, again. At $d = \lambda/2$, $\Delta\varphi = \pi\sin\theta$ stays
      inside $(-\pi, \pi]$: every angle gets its \emph{own} phase step. Space
      them wider and two directions share one reading.}
\end{frame}

% Teaching note (180 s): from one step to a pattern. Three antennas at 30 deg read
% 0, 90, 180 degrees -- a straight line in antenna index. Name "steering vector" as
% nothing more than "the pattern angle theta WOULD produce".
\begin{frame}{M antennas: the phases form a line}
  \vspace{0.3em}
  \centering
  \begin{tikzpicture}[scale=0.9,transform shape]
    % axes
    \draw[-{Latex[length=2mm]},softgray] (0,0) -- (6.6,0)
      node[right,font=\scriptsize]{$m$};
    \draw[-{Latex[length=2mm]},softgray] (0,0) -- (0,2.9)
      node[rotate=90,font=\scriptsize,anchor=south east,yshift=2pt]{phase $\varphi_m$};
    \foreach \m/\x in {1/1.0, 2/2.6, 3/4.2, 4/5.8}{
      \draw[softgray!60] (\x,-0.07) -- (\x,0.07);
      \node[font=\scriptsize,text=softgray] at (\x,-0.33) {\m};}
    % the line through measured points
    \draw[ranblue,thick] (0.6,0.25) -- (6.1,2.6);
    \foreach \x/\y in {1.0/0.42, 2.6/1.10, 4.2/1.78, 5.8/2.47}{
      \fill[ranblue] (\x,\y) circle (0.085);}
    \node[font=\scriptsize\bfseries,text=ranblue,rotate=23] at (3.6,1.95)
      {slope $= \pi\sin\theta$ per antenna};
  \end{tikzpicture}

  \vspace{0.4em}
  {\small
   \begin{itemize}\itemsep3pt
     \item $\varphi_m = m \cdot \pi\sin\theta$: the direction sets the
       \textbf{slope} across the array.
     \item The pattern a candidate angle \emph{would} produce is its
       \textbf{steering vector} $a(\theta)$.
     \item Finding $\theta$ = finding which template line fits the measured
       column of $H$.
   \end{itemize}\par}

  \glossline{Boresight ($\theta = 0$): slope 0 --- all antennas in step.
    Endfire ($\theta = 90^\circ$): half a turn each.}
\end{frame}

% Teaching note (180 s): the algorithm is three lines of code; say so. Sweep theta,
% correlate, plot. Emphasize "after the fact, in software" -- no moving parts, the
% beam is a for-loop. This is the frame to connect back to beamforming in Class 8.
\begin{frame}{Scanning every angle in software}
  \vspace{0.2em}
  {\small For each candidate $\theta$ from $-90^\circ$ to $+90^\circ$:
    how well does $a(\theta)$ match the measured antenna column?\par}

  \vspace{0.4em}
  \centering
  \begin{tikzpicture}[scale=0.9,transform shape]
    \draw[-{Latex[length=2mm]},softgray] (0,0) -- (8.6,0)
      node[below left,font=\scriptsize]{$\theta$};
    \draw[-{Latex[length=2mm]},softgray] (0,0) -- (0,2.5)
      node[rotate=90,font=\scriptsize,anchor=south east,yshift=2pt]{match $P(\theta)$};
    \foreach \t/\x in {-90/0.4, -45/2.3, 0/4.2, 45/6.1, 90/8.0}{
      \draw[softgray!60] (\x,-0.07) -- (\x,0.07);
      \node[font=\scriptsize,text=softgray] at (\x,-0.33) {$\t^\circ$};}
    \draw[ranblue,very thick] plot[smooth] coordinates
      {(0.4,0.35) (1.5,0.45) (2.9,0.38) (4.2,0.55) (4.9,0.8) (5.35,2.1)
       (5.5,2.25) (5.65,2.1) (6.1,0.75) (7.0,0.42) (8.0,0.33)};
    \draw[softgreen,dashed,thick] (5.5,0) -- (5.5,2.25);
    \node[font=\scriptsize\bfseries,text=softgreen] at (6.6,2.3)
      {peak at $+30^\circ$};
  \end{tikzpicture}

  \vspace{0.45em}
  \takeaway{This is a phased-array sweep with no moving parts:\\
    the beam is a \texttt{for}-loop over $\theta$, run \emph{after} reception.}

  \source{https://www.usenix.org/conference/nsdi13/technical-sessions/presentation/xiong}{Xiong and Jamieson, ArrayTrack, NSDI 2013}
\end{frame}

% Teaching note (180 s): break the clean picture. The wall echo is a real arrival
% with its own slope, so the spectrum honestly reports two directions. The question
% "which peak is the phone?" has NO answer on this axis -- leave it hanging; it is
% the cliffhanger that justifies the ToF section and SpotFi.
\begin{frame}{Multipath grows extra peaks}
  \vspace{0.2em}
  \centering
  \begin{tikzpicture}[scale=0.9,transform shape]
    \draw[-{Latex[length=2mm]},softgray] (0,0) -- (8.6,0)
      node[below left,font=\scriptsize]{$\theta$};
    \draw[-{Latex[length=2mm]},softgray] (0,0) -- (0,2.5);
    \foreach \t/\x in {-90/0.4, -45/2.3, 0/4.2, 45/6.1, 90/8.0}{
      \draw[softgray!60] (\x,-0.07) -- (\x,0.07);
      \node[font=\scriptsize,text=softgray] at (\x,-0.33) {$\t^\circ$};}
    \draw[ranblue,very thick] plot[smooth] coordinates
      {(0.4,0.35) (1.6,0.5) (2.2,1.1) (2.5,1.75) (2.8,1.1) (3.4,0.5)
       (4.2,0.5) (4.9,0.8) (5.35,2.1) (5.5,2.25) (5.65,2.1) (6.1,0.75)
       (7.0,0.42) (8.0,0.33)};
    \node[font=\Large\bfseries,text=ranpurple] at (2.5,2.1) {?};
    \node[font=\Large\bfseries,text=softgreen] at (5.5,2.55) {?};
    \node[font=\scriptsize,text=ranpurple] at (2.5,-0.75) {wall echo, $-40^\circ$};
    \node[font=\scriptsize,text=softgreen] at (5.5,-0.75) {phone, $+30^\circ$};
  \end{tikzpicture}

  \vspace{0.4em}
  {\small
   \begin{itemize}\itemsep3pt
     \item Every arrival is real: the echo has its own slope, its own peak.
     \item On this axis alone, \textbf{nothing says which peak is the phone}.
     \item ArrayTrack's brute-force help: more antennas $\to$ sharper peaks
       (16 antennas $\to$ 23 cm median).
   \end{itemize}\par}

  \vspace{0.3em}
  \ask{We need a second ruler that tells the two peaks apart.\\
    Hold that thought --- first, let us sharpen this one.}
\end{frame}

% Teaching note (210 s): MUSIC without eigen-anxiety. The trick to say out loud:
% instead of asking "where IS signal?", ask "which directions are provably EMPTY?",
% then plot 1/emptiness. Hands-on versions: Princeton lab, Tsinghua tutorial code.
\begin{frame}{MUSIC, in one slide}
  \vspace{0.3em}
  {\small
   \begin{itemize}\itemsep4pt
     \item Collect many packets $\to$ average $\to$ a covariance of the antenna
       readings.
     \item Its eigenvectors split cleanly: a few span the \textbf{signal}
       (one per path), the rest span pure \textbf{noise}.
     \item A true arrival angle's steering vector is \emph{orthogonal} to the
       noise part.
     \item So plot $1/(\text{projection onto noise})$: division by almost-zero
       $\to$ needle-sharp spikes.
   \end{itemize}\par}

  \vspace{0.3em}
  \centering
  \begin{tikzpicture}[scale=0.85,transform shape]
    \draw[-{Latex[length=2mm]},softgray] (0,0) -- (8.4,0)
      node[below left,font=\scriptsize]{$\theta$};
    \draw[softgray!80,thick] plot[smooth] coordinates
      {(0.4,0.3) (2.0,0.9) (2.5,1.05) (3.0,0.9) (4.4,0.55) (5.1,1.2)
       (5.5,1.35) (5.9,1.2) (7.0,0.45) (8.0,0.3)};
    \draw[ranblue,very thick] plot[smooth] coordinates
      {(0.4,0.15) (2.3,0.2) (2.5,1.9) (2.7,0.2) (5.2,0.25) (5.5,2.2)
       (5.8,0.25) (8.0,0.15)};
    \node[font=\scriptsize,text=softgray] at (7.3,0.95) {plain scan};
    \node[font=\scriptsize\bfseries,text=ranblue] at (6.75,2.15) {MUSIC};
  \end{tikzpicture}

  \source{https://www.cs.princeton.edu/courses/archive/spring18/cos463/labs/lab5_preview.html}{Schmidt, IEEE Trans.\ AP 1986 \quad$\cdot$\quad Have a try: Princeton COS463 MUSIC lab; Tsinghua tutorial \texttt{naive\_aoa}}
\end{frame}
